\section*{Capítulo \ref{ch:g10:heart-lungs-effort} --- El corazón y los pulmones durante el esfuerzo}

\begin{solution}{exo:g10:heart-lungs-effort:1}
$0.6 \times 14 = \qty{8.4}{L/min}$; $2.2 \times 35 = \qty{77}{L/min}$.
\end{solution}

\begin{solution}{exo:g10:heart-lungs-effort:2}
El volumen de sangre que un lado del corazón bombea por minuto:
$Q = f \times V_s$. $65 \times 0.075 \approx \qty{4.9}{L/min}$.
\end{solution}

\begin{solution}{exo:g10:heart-lungs-effort:3}
Aurícula derecha, ventrículo derecho, arteria pulmonar, pulmones
(circulación pulmonar), venas pulmonares, aurícula izquierda,
ventrículo izquierdo, aorta, órganos (circulación sistémica), venas,
aurícula derecha.
\end{solution}

\begin{solution}{exo:g10:heart-lungs-effort:4}
Volumen corriente \qty{0.5}{L}; capacidad vital \qty{4.5}{L}; volumen
residual \qty{1.2}{L}.
\end{solution}

\begin{solution}{exo:g10:heart-lungs-effort:5}
Unos $220 - 16 = 204$ y $220 - 60 = 160$ latidos por minuto.
\end{solution}

\begin{solution}{exo:g10:heart-lungs-effort:6}
A \qty{100}{W}: $120 \times 0.105 \approx \qty{12.6}{L/min}$; a
\qty{300}{W}: $198 \times 0.115 \approx \qty{22.8}{L/min}$; un factor
de 1,8.
\end{solution}

\begin{solution}{exo:g10:heart-lungs-effort:7}
Reposo: \qty{2.5}{L/min}; ejercicio: $0.04 \times 24 \approx
\qty{1.0}{L/min}$. La proporción cae doce veces, y el flujo absoluto
solo por un factor de 2,6, porque el gasto mismo ha crecido.
\end{solution}

\begin{solution}{exo:g10:heart-lungs-effort:8}
$18 \times 0.14 = \qty{2.5}{L/min}$: unas tres cuartas partes de los
\qty{3.45}{L/min} de un estudiante entrenado --- cerca de ese valor
para una persona más pequeña.
\end{solution}

\begin{solution}{exo:g10:heart-lungs-effort:9}
El cerebro no tolera ni segundos sin oxígeno y glucosa, así que su
aporte está protegido. El tubo digestivo y los riñones pueden suspender
la digestión y la filtración durante una hora sin daño, y su parte se
presta a los músculos.
\end{solution}

\begin{solution}{exo:g10:heart-lungs-effort:10}
$V_s = 5000/42 \approx \qty{120}{mL}$, frente a unos \qty{70}{mL}: su
corazón agrandado entrega el gasto de reposo con muchos menos latidos.
\end{solution}

\begin{solution}{exo:g10:heart-lungs-effort:11}
La orden de moverse del encéfalo, y su anticipación del esfuerzo, se
envían a los centros cardíacos y a las glándulas suprarrenales antes de
que los músculos actúen. Un corredor que dependiera solo de la
detección del dióxido de carbono arrancaría con el gasto de reposo,
correría el primer minuto con la \omterm{def:g10:cell-metabolism:fermentation}{fermentación} y se vería obligado a
aflojar por el ácido láctico antes de que la entrega lo alcanzara.
\end{solution}

\begin{solution}{exo:g10:heart-lungs-effort:12}
$195 \times 0.110 \times 0.150 \approx \qty{3.2}{L/min}$ y
$195 \times 0.160 \times 0.150 \approx \qty{4.7}{L/min}$. El volumen
sistólico es lo que agranda el entrenamiento: el segundo sujeto es el
entrenado en resistencia.
\end{solution}

\begin{solution}{exo:g10:heart-lungs-effort:13}
Oxígeno inspirado: $0.21 \times 100 = \qty{21}{L/min}$; usado:
\qty{3.5}{L/min}, es decir un 17\%. El resto vuelve a espirarse --- el
aire espirado todavía tiene alrededor de un 16\% de oxígeno, y por eso
funciona la respiración boca a boca.
\end{solution}

\begin{solution}{exo:g10:heart-lungs-effort:14}
Gasto de como mucho unos $110 \times 0.12 \approx \qty{13}{L/min}$ en
lugar de 24: el $\dot V\!\mathrm{O_2}$máx cae casi a la mitad, hasta
unos \qty{2}{L/min}. El paciente puede caminar y pedalear suave, pero
no sostener ningún esfuerzo intenso.
\end{solution}

\begin{solution}{exo:g10:heart-lungs-effort:15}
La entrega por litro baja una cuarta parte y la extracción no cambia en
proporción: el $\dot V\!\mathrm{O_2}$máx cae un 25\%. En las semanas
siguientes el cuerpo fabrica más glóbulos rojos, y el oxígeno
transportado por litro vuelve a acercarse a \qty{200}{mL}.
\end{solution}

\begin{solution}{pb:g10:heart-lungs-effort:1}
\textbf{1.} $62 \times 0.072 \approx 4.5$; $105 \times 0.100 = 10.5$;
$140 \times 0.116 \approx 16.2$; $175 \times 0.120 = 21.0$;
$198 \times 0.120 \approx \qty{23.8}{L/min}$.

\textbf{2.} $23.8/4.5 \approx 5.3$: la frecuencia aporta $198/62
\approx 3.2$, el volumen sistólico $120/72 \approx 1.7$, y $3.2 \times
1.7 \approx 5.3$.

\textbf{3.} Por encima de unos \qty{225}{W} el volumen sistólico se
queda en \qty{120}{mL}; solo la frecuencia cardíaca sigue subiendo el
gasto.

\textbf{4.} $220 - 17 = 203$: sus 198 caen dentro de la precisión de la
regla.

\textbf{5.} $5/4.5 \approx \qty{1.1}{min}$, unos \qty{67}{s}; a
\qty{300}{W}, $5/23.8 \approx \qty{0.21}{min}$, unos \qty{13}{s}.

\textbf{6.} $0.5 \times 13 = 6.5$; $1.2 \times 18 = 21.6$; $1.8 \times
24 = 43.2$; $2.3 \times 32 = 73.6$; $2.6 \times 44 \approx
\qty{114}{L/min}$.

\textbf{7.} $114/6.5 \approx 18$: el volumen corriente subió 5,2 veces
y la frecuencia 3,4; el volumen aporta más.

\textbf{8.} $2.6/5.0 = 52\%$ de la capacidad vital por respiración.

\textbf{9.} No: la \omterm{def:g10:heart-lungs-effort:ventilation}{ventilación} tiene una reserva amplia y la sangre que
sale de los pulmones sigue saturada. El límite es el volumen sistólico
del corazón, que fija cuánto oxígeno entrega cada latido.

\textbf{10.} $0.7 \times 6.5 \approx \qty{4.6}{L/min}$ y $0.7 \times
114 \approx \qty{80}{L/min}$.

\textbf{11.} $4.5 \times 0.052 \approx 0.23$; $10.5 \times 0.090
\approx 0.95$; $16.2 \times 0.120 \approx 1.94$; $21.0 \times 0.145
\approx 3.05$; $23.8 \times 0.155 \approx \qty{3.7}{L/min}$.

\textbf{12.} $3700/68 \approx \qty{54}{mL/min}$ por kilogramo. La
frecuencia ha llegado a su máximo y el aumento de 225 a \qty{300}{W} es
pequeño: este es su $\dot V\!\mathrm{O_2}$máx, el de un estudiante en
forma y entrenado.

\textbf{13.} $3.7/0.23 \approx 16$: el factor 5,3 del corazón
multiplicado por el $155/52 \approx 3.0$ de la extracción.

\textbf{14.} $3.7 \times 20 = \qty{74}{kJ}$ por minuto, unos
\qty{1230}{W}; en reposo, $0.23 \times 20/60 \approx \qty{77}{W}$; por
encima del reposo, unos \qty{1150}{W}, así que el rendimiento es
$300/1150 \approx
26\%$.

\textbf{15.} Aproximadamente una recta de \num{0.23} a
\qty{3.05}{L/min} entre 0 y \qty{225}{W}, y después una curvatura: los
últimos \qty{75}{W} solo añaden \qty{0.65}{L/min}. El techo está en
unos \qty{3.7}{L/min}.

\textbf{16.} La anticipación: la expectativa del esfuerzo que tiene el
encéfalo se transmite a los centros cardíacos y a las glándulas
suprarrenales, y sube la frecuencia antes de todo movimiento.

\textbf{17.} Los \omterm{def:g10:heart-lungs-effort:ventilation}{alvéolos} oxigenan la sangre por completo incluso con
la \omterm{def:g10:heart-lungs-effort:ventilation}{ventilación} y el gasto máximos; los pulmones no son el límite. El
límite es la cantidad de sangre que el corazón puede enviar por minuto.

\textbf{18.} $20/23.8 \approx 84\%$. El tubo digestivo y los riñones,
que se llevaban la mitad del gasto de reposo, reciben ahora
\qty{1}{L/min} aproximadamente; sus arterias se han estrechado y el
flujo se ha desviado a los músculos.

\textbf{19.} $198 \times 0.140 \times 0.155 \approx \qty{4.3}{L/min}$,
unos \qty{63}{mL/min} por kilogramo.

\textbf{20.} El \omterm{def:g10:heart-lungs-effort:output}{gasto cardíaco} multiplicado por 5,3, la extracción de
oxígeno por 3 y la \omterm{def:g10:heart-lungs-effort:ventilation}{ventilación} por 18; el entrenamiento cambió el
volumen sistólico del corazón, y con él el techo.
\end{solution}
